\section{The obstruction bound at all Taylor--Wiles levels}

Fix the source's arithmetic Taylor--Wiles data $(T,\mathcal P)$,
a normalized point $x$, its package $k_x$ with $g$ generators,
and its ordered auxiliary roots. Write $A=k[[t]]$, $q=|k|$,
$X_e=A/(t^e)$, $a=x.a$ and $b=n-a$. The odd matrices have the
two off-diagonal blocks of sizes $a\times b$ and $b\times a$.
The problem asks for one exponent $C$ for all truncations $m$.
The threshold level may depend on $m$. The coefficient row must have
a unit entry even when every prescribed obstruction is torsion.

% BEGIN U70O1
\paragraph{U70O1. The global and local duality inputs at every level.}
For any arithmetic level $V$ use its group over $L^+$ unramified
outside $V.S'$, with the original local groups at $v\mid\ell$
and $u\in S_t$. Put $d=[L^+:\mathbb Q]$ and
$d_v=[L_{w(v)}:\mathbb Q_\ell]$ at $v\mid\ell$.
All selected places split in $L/L^+$. Thus these completions and
decomposition groups identify with their counterparts over $L^+$,
and $\sum d_v=d$ by the local degree formula. An $\ell$-killed
finite coefficient module, over any auxiliary ring, has the same
additive continuous cochains and differentials when regarded as an
$\mathbb F_\ell$-module. The finiteness, Euler characteristics and
degree bounds in the original $\mathrm{CohP}$ therefore apply to
these cochains. In particular $H^1,H^2$ are finite locally and
globally. Global and local $H^2$ carry coefficient surjections
onto surjections by the exact sequence and $H^3=0$ for the kernel.
Globally one uses $\ell\ne2$ and the inclusion of every
$\ell$-adic place in $V.S'$. The real term in the global Euler
formula is the product of the fixed-space orders at the $d$ real
involutions. The local Euler exponent is $d_v$ at $v\mid\ell$
and zero at $u\in S_t$.

For the exact subgroup $J_u$ in the source, inertia modulo $J_u$
embeds in the product of the residual character images in the
multiplicative group of the residue field of $O$. Its order is prime
to $\ell$. For an $\ell$-killed module $M$ on which $J_u$ acts
trivially, a cocycle vanishing on $J_u$ restricts on this finite
inertia quotient to a coboundary by averaging. The restrictions
have $|M|/|M^I|$ possibilities and all are attained by global
coboundaries on the local group. After removing a coboundary,
the cocycle factors through the procyclic residue group on $M^I$.
Its value at Frobenius is arbitrary in $M^I$: the generator and
its cocycle value map to an element of a finite semidirect product,
which determines a continuous map from $\widehat{\mathbb Z}$.
Hence there are exactly $|M|$ cocycles vanishing on $J_u$, and the
unramified $H^1$ has order $|H^0(M)|$. For a surjection of such
modules the kernel satisfies the same count. The order formula
then shows that every target cocycle vanishing on $J_u$ lifts to
one vanishing on $J_u$ upstairs.

For finite data consisting of a global module $M$ and surjective
equivariant local quotients $Q_w$, put
\[
 K_2=\operatorname{coker}\!\left(H^2(M)\longrightarrow
         \prod_w H^2(Q_w)\right),\qquad
 M^*=\operatorname{Hom}_{\mathbb F_\ell}(M,\mathbb F_\ell(1)).
\]
The normalized Poitou--Tate sequence in $\mathrm{CohP}$ bounds
$|K_2|$ by $|H^0(M^*)|$. Indeed take all local places in $V.S'$
first, identify the cyclotomic twist using a primitive $\ell$th
root, and apply exactness at the product of local $H^2$. Its
cokernel embeds in the dual of $H^0(M^*)$, of the same order.
Restriction to the specified factors and passage to $Q_w$ are
surjective by local right exactness, so their cokernel is a
quotient of the first. This gives each field of $\mathrm{CohAx}$
for $V$. At level $N$, the given map from the arithmetic group
for $\mathcal P.V(N)$ to $T.G_N$ is a continuous bijection of
profinite groups, by its injectivity and the surjectivity of
the frame map. Transport continuous cochains in both directions
along this homeomorphism. Its compatibility with the local maps,
the sign and the cyclotomic character transports the same formulas.
They hold both at the base and at every required $N$.
% END U70O1

% BEGIN U70O2
\paragraph{U70O2. The obstruction group with local conditions.}
For a projection of frames to the base, use $M=\mathrm{Od}(X_e)$
with conjugation by $r$ on the index-two subgroup and action
$-{}^T$ at the specified involution. The form identity gives a
group action and continuity follows modulo $t^e$. At $v\mid\ell$
let $Q_v$ be the quotient by the odd matrices upper triangular
in the given ordered flag basis. At $u\in S_t$ use $Q_u=M$.
The latter is trivial on $J_u$. These data define the native
odd coefficient datum. Write $J_v$ for the entire local group at
$v\mid\ell$, and set
\[
 S=\ker\!\left(H^1(M)\to\prod_w H^1(Q_w)/H^1_{J_w}(Q_w)\right)
\]
where $H^1_{J_w}$ denotes the kernel of restriction to $J_w$.
The obstruction group $H_e$ is represented by pairs $(z,(y_w))$,
where $z$ is a continuous global $2$-cocycle and
$d y_w=\pi_w\operatorname{res}_w z$. A pair is zero when
$z=d x$ and each $y_w-\pi_w\operatorname{res}_w x$ is a
$1$-cocycle restricting to a coboundary on $J_w$.
All these cochains have coefficients killed by $t^e$.

Put $P=|X_e|^{ab}$. Then $|M|=P^2$ and $|Q_v|=P$.
In the flag basis each column lies in one of the two blocks.
For each cross-block pair of indices exactly one of its two
ordered entries is strictly lower triangular. At any real
involution the action interchanges the two off-diagonal blocks
by inverse isomorphisms, through negative transpose followed by
block-diagonal conjugation. Its fixed subgroup has order $P$.
These statements include $b=0$.

To count, the forgetful map $H_e\to H^2(M)$ lands in the
kernel of localization, and pairs with $z$ a coboundary come
from the product of the local $H^1(Q_w)/H^1_{J_w}(Q_w)$.
Its map to $H_e$ kills restrictions of global $H^1(M)$,
whose kernel in the product is $S$. Thus
\[
 |H_e|\le
 \frac{|S|\,|H^2(M)|\,|K_2|}{|H^1(M)|}
 \prod_w \frac{|H^1(Q_w)|}{|H^1_{J_w}(Q_w)|\,|H^2(Q_w)|}.
\]
All groups here are finite by U70O1 and this kernel argument.
Global Euler gives $|H^1(M)|=|H^0(M)|\,|H^2(M)|P^d$.
The local Euler quotients contribute at most
$|H^0(Q_v)|P^{d_v}$ at $v\mid\ell$ and one at $u$ by
U70O1. Cancel $P^d$ and use its bound for $K_2$. We obtain
\[
          |H_e|\le |S|\,\prod_{v\mid\ell}|H^0(Q_v)|\,
                     |H^0(M^*)|.
\]
% END U70O2

% BEGIN U70O3
\paragraph{U70O3. The uniform base-level Selmer count.}
There is $C_0$ depending on the fixed base data and package such
that $|S_{0,e}|\le q^{ge+C_0}$ for all $e$. To see the uniformity,
record a matrix cocycle on the index-two subgroup together with
one correcting odd matrix at each ordinary place. These variables
are constrained by homogeneous integral linear equations: the
cocycle and polarization laws, block support, zero on the
prescribed inertia, and the strict lower and torsion-diagonal
entries of the locally corrected cocycle. Each equation is a
finite sum of evaluations. A solution modulo $t^e$ is locally
constant. Lifting each value by its polynomial of degree $<e$
gives continuous integral variables satisfying the equations
modulo $t^e$. Conversely reducing such variables gives a solution.

The space of integral variable-functions continuous modulo every
power, with the finitely many local matrices, is complete,
separated and $t$-torsion-free: take compatible pointwise limits
and use their locally constant reductions. The target product of
the equation modules is separated and torsion-free. Solutions at
the first truncation are finite. Indeed a polarized matrix cocycle
extends to the full group with value zero at $c$ using its
transpose identity. The finite $H^1$ from U70O1 and the finite
space of coboundaries then bound the number of cocycles. There
are only finitely many possible local matrices.

In the exact integral kernel take the submodule of tuples whose
global function lies in the span of the $\kappa_j$ plus the
original skew odd coboundaries. The $\mathrm{Kap.gen}$ field
puts a power of $t$ times each exact tuple in that submodule.
CFO3 now bounds the index of its reductions inside the truncated
solutions by $q^{C_0}$, uniformly in $e$. Projection away from
the local matrices can only decrease this index. The projected
reductions have cardinal at most
$q^{ge}|\mathrm{Bd}_e|$, where $\mathrm{Bd}_e$ is the group
of skew odd coboundaries at truncation $e$.

Finally a full-group class in $S_{0,e}$ has a representative
vanishing at $c$, since $2$ is invertible. On the index-two
subgroup it then satisfies the transpose condition. Its local
class gives an upper triangular correction at $v\mid\ell$.
In the typed flag basis an odd matrix has zero diagonal, so the
torsion-diagonal condition follows. At $u$ the local action on
$J_u$ is trivial, so vanishing in cohomology there is vanishing
of the cocycle. Two normalized representatives of the same class
differ by a coboundary with primitive fixed by $c$, exactly a
skew odd primitive. Their restrictions differ by the same group
$\mathrm{Bd}_e$; the restriction determines a normalized cocycle
on both cosets. Dividing the preceding cardinal bound by
$|\mathrm{Bd}_e|$ proves the assertion.
% END U70O3

% BEGIN U70O4
\paragraph{U70O4. The change of level costs a bounded number of classes.}
Write $p_N:T.G_N\to G_0$ for the projection of full groups and
$\tau_{Nj},\varphi_{Nj},\gamma_j$ for the Taylor--Wiles inertia,
Frobenius and chosen base elements. For each $m$, at all sufficiently
large $N$, the matrices of $p_N(\varphi_{Nj})$ and $\gamma_j$
agree modulo $t^m$ for every $j$. Indeed the congruence kernel
in the index-two subgroup contains the relative pro-$\ell$
residual by the given hypothesis. Its normal core in the full
group is open and still contains that normal residual. Apply
the given Frobenius convergence, simultaneously for the finite
list. This equality persists at every $e\le m$.

The action on $M$ at level $N$ is pulled back from the base.
For any continuous cocycle into $M$, the combined action and
cocycle give a homomorphism to a finite affine group.
On $p_N^{-1}(T.H)$ its linear image is an $\ell$-group, since
the base linear representation there factors through the finite
quotients of $T.H$ of $\ell$-power order by the relative residual
hypothesis. Its translation image in the kernel of the linear action
is also an $\ell$-group. Thus the combined map kills the level's
relative pro-$\ell$ residual. Since $p_N(\tau_{Nj})=1$,
evaluation on these elements factors through cohomology.
The tame relation says that conjugation by $\varphi_{Nj}^{-1}$
takes $\tau_{Nj}$ to its $Nq_{Nj}$th power modulo that
residual. For $N\ge1$, $Nq_{Nj}\equiv1\pmod\ell$.
The cocycle law thus puts its value at $\tau_{Nj}$ in the odd
centralizer of $r(\gamma_j)$ modulo $t^e$.

Over $E=\operatorname{Frac} A$ this odd centralizer is zero.
The eigenvalues of $r(\gamma_j)$ are all distinct by the given
$\mathrm{Eig}$, so its two blocks have disjoint eigenvalues.
The commutator map on the finite free odd matrix module has
an invertible matrix over $E$ whenever the module is nonzero.
A common denominator for its inverse bounds the size of every
truncated kernel by $q^{C_j}$ independently of $e$ (each of its
coordinates is killed by that denominator). For a zero module
take $C_j=0$. If all evaluations vanish, the affine map kills
the normal closure of the $\tau_{Nj}$ as well as the residual.
Their closed join contains $\ker p_N$ by the Taylor--Wiles
kernel axiom, so the cocycle descends continuously. The local
maps and local quotient modules commute with $p_N$, which
preserves the Selmer conditions in descent. Conversely its
inflation is the original class. Counting the fibers of the
joint evaluation gives, for $C_T=\sum C_j$ and simultaneously
for $0\le e\le m$ at these levels,
\[
                  |S_{N,e}|\le |S_{0,e}| q^{C_T}.
\]
% END U70O4

% BEGIN U70O5
\paragraph{U70O5. The local invariant and global dual factors.}
There is one constant $C_h$ with $|H^0(Q_v)|\le q^{C_h}$
for all $v,e,N$. In the ordered flag basis $Q_v$ is freely
coordinatized over $X_e$ by the strictly lower cross-block
entries. These coordinates commute with reduction in $e$.
Over $E$ the triangular action on the strictly lower quotient
has a filtration with one-dimensional quotients having characters
$\psi_i/\psi_j$, $i>j$ of different types. For example,
conjugation of $E_{ij}$ by an upper triangular matrix contributes
only entries $(k,l)$ with $k\le i$, $l\ge j$; the coefficient
at $(i,j)$ is $\psi_i/\psi_j$. The other strictly lower
contributions have smaller $k-l$. None of these characters is
trivial by $x.N1$, hence no fixed vector survives over $E$.
The simultaneous linear equations defining fixed vectors thus
have finitely many rows of full rank over $E$. An invertible
minor clears denominators and bounds every truncated kernel,
as in U70O4. There are finitely many $v$.

There is also $C_d$ with $|H^0(M^*)|\le q^{C_d}$ for all $e,N$.
An invariant dual functional $f$ is additive, with the prescribed
cyclotomic equivariance. On the full cyclotomic kernel this means
$f(r(h)Y)=f(Yr(h))$ for every odd $Y$, after transporting from
the level along its surjective map respecting the cyclotomic
character. The latter means lifts of cyclotomic-kernel elements
are still in that kernel. Replacing $Y$ by scalar multiples and
adding proves the same sliding identity for the $A$-span of these
matrices, despite $f$ being only $\mathbb F_\ell$-linear. By
$x.\mathrm{span}$ this span contains $t^sP_1$ for a fixed $s$,
where $P_1$ is the projector on the first block. Sliding $t^sP_1$
across $Y$ in each pure corner separately forces $f(t^sY)=0$.
Thus $f$ factors through the finite additive space of odd matrices
modulo $t^{\min(e,s)}$, of bounded cardinality and bounded
$\mathbb F_\ell$-dimension. Counting all functionals on that
space gives the asserted bound.

At every level of U70O4 we now have, uniformly for $e\le m$,
\[
 \begin{split}
 |H_{N,e}|&\le q^{ge+C_0+C_T+|V_\ell|C_h+C_d},\\
 |K_{2,N,e}|&\le q^{C_d}.
 \end{split}
\]
Here $V_\ell$ denotes the fixed set of places above $\ell$ and
the second line uses U70O1 directly.
% END U70O5

% BEGIN U70O6
\paragraph{U70O6. Bounds for every quotient of the obstruction group.}
Put $H=H_{N,m}$. For each $0\le e\le m$ there are two short
exact sequences of coefficient modules
\[
 \begin{split}
 0\to X_{m-e}\xrightarrow{t^e}X_m\to X_e\to0,\\
 0\to X_e\xrightarrow{t^{m-e}}X_m\to X_{m-e}\to0.
 \end{split}
\]
They are exact on odd global modules and on local quotients,
by their free coordinate descriptions. Write $f_*,r_*$ for the
maps on obstruction groups induced by the first inclusion and
the second surjection. Their composite $f_*r_*$ on $H$ is
multiplication by $t^e$, on both coordinates of its defining pairs.

For the first sequence, the kernel of $H\to H_{N,e}$ lies in
$\operatorname{im} f_*$. Indeed if a pair reduces to zero, lift
its global primitive, using any section between the finite modules,
to make its global cocycle kernel-valued. The reduced remaining
local cochains are cocycles which are coboundaries on $J_w$.
They lift to such cocycles on the middle local modules: lift the
local coboundary primitive first, and then lift the cocycle zero
on $J_w$ using U70O1 at $u$. At $v\mid\ell$ the remaining
cocycle is identically zero since $J_v$ is the entire group.
Subtracting these allowed terms gives kernel-valued local cochains
and hence a lifting pair from the first module. All choices of
sections for cochains preserve continuity because the modules are finite.

For the second sequence the cokernel of $r_*$ has cardinal at
most $|K_{2,N,e}|$. For each target pair, lift its global
$2$-cocycle to a global $2$-cocycle, first in cohomology by
U70O1 and then exactly by lifting a correcting $1$-cochain.
Lift its local cochains using set sections. At each place the
difference between the localized global cocycle and the differential
of the chosen local cochain lies in the coefficient kernel
$Q_{w,e}$ and is a $2$-cocycle. Take its class in
$\prod_wH^2(Q_{w,e})$ modulo global localization.
If two choices of target pairs give the same resulting element in
$K_{2,N,e}$, subtract them. A global kernel-valued $2$-cocycle
now accounts for the local defect classes; subtract it and adjust
the lifted local cochains by kernel-valued $1$-cochains. This
makes the difference a lifting pair. Choosing one target pair and
one set of lifts per cokernel coset thus gives an injection into
$K_{2,N,e}$.

Consequently $[H:\operatorname{im} f_*]\le|H_{N,e}|$, while
$[\operatorname{im} f_*:t^e H]\le|K_{2,N,e}|$ by the
factorization through $\operatorname{coker} r_*$. In particular
$H$ is finite, $t^mH=0$, and for
$C=C_0+C_T+|V_\ell|C_h+2C_d$ we have
\[
                       |H/t^eH|\le q^{ge+C}
                       \qquad(0\le e\le m).
\]
All four constants were chosen before $m$ or $N$.
% END U70O6

% BEGIN U70O7
\paragraph{U70O7. The obstruction of an actual array of even cocycles.}
Take any rows $z_{ij}\in\mathrm{SelP}$ at this truncation
and level, $0\le i\le g$, $1\le j\le g$.
With $e_j,f_j,\eta$ as in Q70C2, the source's $W_2$ consists of
$\sum_j f_jP_j+\eta Z$ with $P_j$ even and $Z$ odd.
Under conjugation by $(1+\sum e_j\kappa_j)r$ and negative
transpose at $c$, the $f_j$ coordinates carry the given even
action and the kernel $\eta M$ carries the odd action. These
operations preserve $W_2$, since products of odd and even
matrices are odd. Thus there is an equivariant exact sequence
\[
         0\longrightarrow M\xrightarrow{\,\eta\,}
              W_2\longrightarrow \mathrm{Ev}(X_m)^g\longrightarrow0.
\]
Extend the polarized row to a full-group cocycle by setting its
value at $c$ equal to zero. Lift to a cochain $x_i$ of $W_2$ by
putting its $\eta$ coordinate zero. Then $d x_i=\eta u_i$ with
$u_i$ a continuous global $2$-cocycle.

At $v\mid\ell$ lift the row's even ordinary primitives by the
same section and let $b_{i,v}$ be their coboundary. Modulo the
subspace upper triangular in the deformed flag basis
$w_1=(1-\sum e_jY_{j,v})w$, the difference $x_i-b_{i,v}$
is represented by $\eta M$. Indeed its even coordinates are
upper triangular after the ordinary corrections. Conjugating by
$w_1$ adds to the odd coordinate just commutators between odd
and even matrices in addition to the change by $w$. One can
therefore match the remaining lower cross-block entries by an
odd matrix. Its class in $Q_v$ is independent of the choice
since the kernel is precisely upper triangular odd matrices.
At $u$ the inertia $J_u$ acts trivially on all of $W_2$: both
$r$ and $\kappa_j$ have the required trivial action and zero
cocycle values there, respectively. The row is zero there.
U70O1 lifts its local even cocycle to a $W_2$-cocycle $b_{i,u}$
zero there, giving $x_i-b_{i,u}=\eta y_{i,u}$. Write $y_{i,v}$
for the preceding class at $v$. Because the triangular subspace
is stable and $d b_{i,w}=0$, we have
$d y_{i,w}=\pi_w\operatorname{res}_w u_i$. Thus
$\xi_i=[u_i,(y_{i,w})]\in H$.

If $\sum_i d_i\xi_i=0$ for coefficients $d_i\in A$, take a
global primitive $p$ in the definition of this zero class.
Then $\sum_i d_i x_i-\eta p$ is a cocycle with exactly the
even coordinates $\sum_i d_i z_{ij}$. Its value at $c$ is
$\eta M$-valued and anti-invariant under $c$. Adding the
coboundary of half that value makes it vanish at $c$, giving the
prescribed polarization on restriction. The local primitive in
$Q_w$ for each residual cochain on $J_w$ lifts to $M$.
At $u$ this correction has zero coboundary on $J_u$, giving
the required value zero. At $v$ it joins the sum of the chosen
ordinary primitives, giving upper triangularity in the deformed
flag. The corrected diagonal on $F.\mu_v$ also vanishes:
its $f_j$ entries are exactly the corrected even diagonals,
and its $\eta$ entry in the typed flag basis belongs to an
off-diagonal block and therefore has zero diagonal. The added
global coboundary merely shifts each local primitive.
Continuity holds because the coefficients are finite and all the
cochains are continuous. This constructs the native
$\Psi\in\mathrm{Sel2}$ for the row $d$.
% END U70O7

% BEGIN U70O8
\paragraph{U70O8. A unit coefficient at all truncations and levels.}
Let $H$ be the finite $A$-module of U70O6 and take the $g+1$
elements $\xi_i$ from U70O7. In their span $D$ write the
invariant factors as $A/(t^{b_j})$ in decreasing order.
There are at most $g+1$ by generation. If there are fewer,
the map $A^{g+1}\to D$ has a primitive relation, by Smith
normal form over $A$. If there are $g+1$, the same normal form
gives a primitive vector whose image generates a summand of
length $b_{g+1}$. We claim $b_{g+1}\le C$. Otherwise
$e=C+1\le m$ since $t^mH=0$, and
\[
  q^{(g+1)e}=|D/t^eD|=|D[t^e]|\le |H[t^e]|
              =|H/t^eH|\le q^{ge+C},
\]
an impossibility. The two equalities comparing kernel and quotient
use finiteness for multiplication by $t^e$. If $D=0$ take any
standard basis vector. We have in every case a coefficient vector
$c$ with a unit entry and $t^C\sum c_i\xi_i=0$.
U70O7 with $d_i=t^Cc_i$ supplies a single $\Psi$ with
\[
         (\Psi(\gamma)_{ab}).f_j
                 =\left(t^C\sum_i c_i z_{ij}(\gamma)\right)_{ab}
\]
for every $j,\gamma,a,b$. This includes $m=0$ and $g=0$.

Thus one $C$ works for every $m$, at all sufficiently large
$N$ for that $m$, and for every array at such a level. The
proof used the original $\mathrm{CohP}$, the given arithmetic
places, the relative pro-$\ell$ kernel condition and the ordered
distinct eigenvalues. It proves the original
$\mathrm{TaylorWilesObstructionBound}$ and supplies $B2$ in
N69E4 with the same normalized point and package.
% END U70O8
